\textbf{Phần 1. Trắc nghiệm lựa chọn}
\setcounterpageref{ex}{0}

\begin{ex} %[2D4V21]
Tính giá trị biểu thức: $37 + 19 + 24$.
\choice
{\True $80$}
{$78$}
{$79$}
{$81$}
\loigiai{
$37 + 19 = 56$; $56 + 24 = 80$.
\textbf{Đáp án của câu hỏi là A}
}
\end{ex}

\begin{ex} %[2D4V21]
Một cuộn dây dài $48\text{ m}$, người ta nối thêm một đoạn $25\text{ m}$ rồi lại nối thêm một đoạn $17\text{ m}$. Cuộn dây lúc này dài bao nhiêu mét?
\choice
{$88\text{ m}$}
{$89\text{ m}$}
{\True $90\text{ m}$}
{$91\text{ m}$}
\loigiai{
$48 + 25 + 17 = 48 + (25 + 17) = 48 + 42 = 90\text{ (m)}$.
\textbf{Đáp án của câu hỏi là C}
}
\end{ex}

\begin{ex} %[2D4V21]
Tìm số tự nhiên $y$ biết: $y - 28 = 47$.
\choice
{$74$}
{\True $75$}
{$76$}
{$65$}
\loigiai{
$y = 47 + 28 = 75$.
\textbf{Đáp án của câu hỏi là B}
}
\end{ex}

\begin{ex} %[2D4V21]
Một cửa hàng bán được $39$ hộp bánh trong ngày thứ nhất, ngày thứ hai bán được nhiều hơn ngày thứ nhất $15$ hộp. Cả hai ngày cửa hàng bán được bao nhiêu hộp bánh?
\choice
{$83$ hộp}
{$92$ hộp}
{$94$ hộp}
{\True $93$ hộp}
\loigiai{
Ngày thứ hai bán: $39 + 15 = 54$ (hộp).\nCả hai ngày bán: $39 + 54 = 93$ (hộp).
\textbf{Đáp án của câu hỏi là D}
}
\end{ex}

\begin{ex} %[2D4V21]
Tính chu vi hình tam giác có độ dài các cạnh lần lượt là $27\text{ cm}, 35\text{ cm}, 18\text{ cm}$.
\choice
{\True $80\text{ cm}$}
{$78\text{ cm}$}
{$82\text{ cm}$}
{$79\text{ cm}$}
\loigiai{
Chu vi: $27 + 35 + 18 = 80\text{ (cm)}$.
\textbf{Đáp án của câu hỏi là A}
}
\end{ex}

\begin{ex} %[2D4V21]
Tìm một số biết rằng lấy số đó trừ đi $38$ thì bằng tổng của $26$ và $19$.
\choice
{$81$}
{$82$}
{\True $83$}
{$84$}
\loigiai{
Tổng hai số: $26 + 19 = 45$. Số cần tìm: $45 + 38 = 83$.
\textbf{Đáp án của câu hỏi là C}
}
\end{ex}

\begin{ex} %[2D4V21]
Điền số thích hợp: $46 + 37 = 28 + \dots$
\choice
{$54$}
{\True $55$}
{$56$}
{$53$}
\loigiai{
Vế trái: $46 + 37 = 83$. Ta có: $83 - 28 = 55$. Vậy điền 55.
\textbf{Đáp án của câu hỏi là B}
}
\end{ex}

\begin{ex} %[2D4V21]
Bác nông dân có $48$ con gà và $37$ con vịt. Bác mua thêm $9$ con ngỗng. Hỏi bác có tất cả bao nhiêu con gia cầm?
\choice
{$92$ con}
{$93$ con}
{$95$ con}
{\True $94$ con}
\loigiai{
$48 + 37 + 9 = 85 + 9 = 94$ (con).
\textbf{Đáp án của câu hỏi là D}
}
\end{ex}

\begin{ex} %[2D4V21]
Tổng của số chẵn lớn nhất có hai chữ số mà chữ số hàng chục là $5$ và số lẻ nhỏ nhất có hai chữ số mà chữ số hàng chục là $3$ là
\choice
{\True $89$}
{$87$}
{$90$}
{$88$}
\loigiai{
Số chẵn lớn nhất có 2 chữ số hàng chục là 5 là 58. Số lẻ nhỏ nhất có 2 chữ số hàng chục là 3 là 31.\nTổng: $58 + 31 = 89$.
\textbf{Đáp án của câu hỏi là A}
}
\end{ex}

\begin{ex} %[2D4V21]
Một xe khách xuất phát với $28$ người. Đến bến đỗ có $15$ người lên xe. Xe đi tiếp và đón thêm $14$ người nữa. Trên xe có bao nhiêu hành khách?
\choice
{$55$ người}
{$56$ người}
{\True $57$ người}
{$58$ người}
\loigiai{
$28 + 15 + 14 = 43 + 14 = 57$ (người).
\textbf{Đáp án của câu hỏi là C}
}
\end{ex}

\begin{ex} %[2D4V21]
Hai số có hiệu bằng $27$, số trừ là $46$. Số bị trừ là
\choice
{$71$}
{\True $73$}
{$72$}
{$74$}
\loigiai{
Số bị trừ: $27 + 46 = 73$.
\textbf{Đáp án của câu hỏi là B}
}
\end{ex}

\begin{ex} %[2D4V21]
Điền dấu thích hợp: $38 + 47 \dots 59 + 25$.
\choice
{$<$}
{$=$}
{Không so sánh được}
{\True $>$}
\loigiai{
$38 + 47 = 85$; $59 + 25 = 84$. Vì $85 > 84$ nên điền dấu $>$. 
\textbf{Đáp án của câu hỏi là D}
}
\end{ex}

\textbf{Phần 2. Trắc nghiệm đúng sai}
\setcounterpageref{ex}{0}

\begin{ex} %[2D4V21]
Xét tính đúng sai của các bài toán tổng hợp cộng có nhớ:
\choiceTF
{\True Biểu thức $28 + 34 + 18$ có kết quả bằng $80$.}
{Nếu một số hạng tăng 5 đơn vị và số hạng kia tăng 7 đơn vị thì tổng tăng 10 đơn vị.}
{\True Phép tính $56 + 27$ có chữ số tận cùng là 3.}
{\True Tổng của hai số chẵn luôn luôn là một số chẵn.}
\loigiai{
\begin{itemize}
    \item Ý a) Đúng: $28 + 34 + 18 = 80$.
    \item Ý b) Sai: Tổng tăng thêm: $5 + 7 = 12$ đơn vị (không phải 10).
    \item Ý c) Đúng: $6 + 7 = 13$, chữ số tận cùng là 3. Đúng.
    \item Ý d) Đúng: Đúng.
\end{itemize}
\begin{itemize}
    \item a) \textbf{Đúng}.
    \item b) \textbf{Sai}.
    \item c) \textbf{Đúng}.
    \item d) \textbf{Đúng}.
\end{itemize}
}
\end{ex}

\begin{ex} %[2D4V21]
Một thư viện trường có 38 quyển truyện tranh và 45 quyển truyện cổ tích:
\choiceTF
{\True Cả hai loại truyện có tất cả 83 quyển.}
{\True Truyện cổ tích nhiều hơn truyện tranh 7 quyển.}
{Nhập thêm 5 quyển truyện tranh thì số truyện tranh sẽ nhiều hơn truyện cổ tích.}
{Tổng số sách chia đều được cho 2 học sinh mà không dư.}
\loigiai{
\begin{itemize}
    \item Ý a) Đúng: $38 + 45 = 83$ quyển.
    \item Ý b) Đúng: $45 - 38 = 7$ quyển.
    \item Ý c) Sai: Truyện tranh có $38 + 5 = 43 < 45$. Sai.
    \item Ý d) Sai: 83 là số lẻ nên không chia hết cho 2.
\end{itemize}
\begin{itemize}
    \item a) \textbf{Đúng}.
    \item b) \textbf{Đúng}.
    \item c) \textbf{Sai}.
    \item d) \textbf{Sai}.
\end{itemize}
}
\end{ex}

\textbf{Phần 3. Trả lời ngắn}
\setcounterpageref{ex}{0}

\begin{ex} %[2D4V21]
Tính giá trị: $47 + 38$.
\shortans{85}
\loigiai{
\textbf{Đáp số:} $85$\\
$47 + 38 = 85$.
}
\end{ex}

\begin{ex} %[2D4V21]
Một trang trại có 36 con bò và 49 con dê. Hỏi trang trại có tất cả bao nhiêu con bò và dê?
\shortans{85}
\loigiai{
\textbf{Đáp số:} $85$\\
$36 + 49 = 85$ con.
}
\end{ex}

\begin{ex} %[2D4V21]
Tìm số tự nhiên a biết: $a - 37 = 48$.
\shortans{85}
\loigiai{
\textbf{Đáp số:} $85$\\
$a = 48 + 37 = 85$.
}
\end{ex}

\begin{ex} %[2D4V21]
Một cuộn dây dài 28 m, người ta nối thêm vào cuộn dây 35 m nữa rồi nối tiếp 19 m. Hỏi cuộn dây dài tất cả bao nhiêu mét?
\shortans{82}
\loigiai{
\textbf{Đáp số:} $82$\\
$28 + 35 + 19 = 82\text{ m}$.
}
\end{ex}

\begin{ex} %[2D4V21]
Tìm tổng của số lớn nhất có hai chữ số mà tổng hai chữ số bằng 8 và số 18.
\shortans{98}
\loigiai{
\textbf{Đáp số:} $98$\\
Số lớn nhất có 2 chữ số tổng bằng 8 là 80. Tổng: $80 + 18 = 98$.
}
\end{ex}

\begin{ex} %[2D4V21]
Bạn Nam sưu tầm được 46 con tem, bạn Hùng sưu tầm được 38 con tem. Hỏi cả hai bạn sưu tầm được bao nhiêu con tem?
\shortans{84}
\loigiai{
\textbf{Đáp số:} $84$\\
$46 + 38 = 84$ con tem.
}
\end{ex}
